\documentclass[11pt,a4paper]{coverletter}

\senderblock
  {Radheem Bin Razi}
  {Distributed Systems \& Cloud-Native Engineer}
  {Ilmenau, Germany \textbar\ \href{mailto:sheikh.radheem@gmail.com}{sheikh.radheem@gmail.com} \textbar\ \href{https://radheem.github.io/my-cv}{portfolio}}

\begin{document}

%% ===== ENGLISH =====
\recipient{Delivery Hero}{Hiring Team}{}

\opening{Dear Hiring Team,}

\vspace{8pt}\noindent\textbf{1. Why Delivery Hero?}\par\vspace{3pt}

The hardest part of AI adoption is not building the models. It is turning scattered experimental tools into reliable, everyday workflows that engineers actually trust. That is exactly why I am applying for the Data Engineer II role within the Catalyst for Agentic Transformation program.

\vspace{8pt}\noindent\textbf{2. Why me?}\par\vspace{3pt}

I have spent the last five years building the quiet infrastructure that turns messy inputs into reliable signals. At a stealth startup, I designed an event-driven platform that exposed declarative tools to LLM agents through a Model Context Protocol server. I also built an automated ingestion pipeline that parsed unstructured email bodies and consolidated them into a single discovery queue. Those systems required strict data-quality gates and consistent routing to prevent downstream failures.

Your team needs to architect pipelines that parse highly unstructured developer logs and extract semantic patterns for root cause analysis. I have already engineered Kafka-based pub/sub architectures that fan out per-user metrics to multiple storage backends while maintaining strict ordering. I understand that technical capability only matters when it becomes a daily habit. I enjoy designing the data foundations that let non-technical stakeholders track productivity and see clear ROI from AI investments.

\vspace{8pt}\noindent\textbf{3. Why now?}\par\vspace{3pt}

I am ready to apply this pipeline experience to a large-scale, cross-brand AI initiative. The timing aligns perfectly with my goal to focus on data storytelling and agentic infrastructure at scale. I am available to start full-time immediately and remain open to part-time arrangements during the transition. I am currently based in Germany and can relocate within the country within two to three weeks. I am Blue Card-ready, which means your legal team only needs to issue the contract without any sponsorship overhead.

\closing{Sincerely,}{Radheem Bin Razi}

\clearpage
\selectlanguage{ngerman}

%% ===== DEUTSCH =====
\recipient{Delivery Hero}{Personalabteilung}{}

\opening{Sehr geehrtes Hiring-Team,}

\vspace{8pt}\noindent\textbf{1. Warum Delivery Hero?}\par\vspace{3pt}

Der schwierigste Teil der KI-Adoption liegt nicht im Bau der Modelle. Es geht darum, verstreute experimentelle Tools in zuverlässige, alltägliche Arbeitsabläufe zu überführen, denen Ingenieure tatsächlich vertrauen. Genau aus diesem Grund bewerbe ich mich für die Position als Data Engineer II im Programm „Catalyst for Agentic Transformation“.

\vspace{8pt}\noindent\textbf{2. Warum ich?}\par\vspace{3pt}

In den letzten fünf Jahren habe ich die unsichtbare Infrastruktur aufgebaut, die unstrukturierte Eingaben in zuverlässige Signale verwandelt. Bei einem Stealth-Startup habe ich eine ereignisgesteuerte Plattform entwickelt, die deklarative Tools über einen Model Context Protocol Server für LLM-Agents bereitstellte. Zudem habe ich eine automatisierte Ingestion-Pipeline implementiert, die unstrukturierte E-Mail-Inhalte parste und sie in einer einzigen Discovery Queue konsolidierte. Diese Systeme erforderten strenge Data-Quality-Gates und ein konsistentes Routing, um Ausfälle in nachgelagerten Prozessen zu verhindern.

Ihr Team benötigt Pipelines, die hochgradig unstrukturierte Developer Logs parsen und semantische Muster für die Root Cause Analysis extrahieren. Ich habe bereits Kafka-basierte pub/sub-Architekturen entwickelt, die nutzerspezifische Metriken auf mehrere Storage Backends verteilen und dabei eine strikte Reihenfolge gewährleisten. Mir ist bewusst, dass technische Fähigkeiten erst dann relevant sind, wenn sie zur täglichen Routine werden. Ich gestalte gerne die Datenfundamente, die es nicht-technischen Stakeholdern ermöglichen, die Produktivität zu verfolgen und einen klaren ROI aus KI-Investitionen zu erkennen.

\vspace{8pt}\noindent\textbf{3. Warum jetzt?}\par\vspace{3pt}

Ich bin bereit, diese Pipeline-Erfahrung in eine groß angelegte, markenübergreifende KI-Initiative einzubringen. Der Zeitpunkt passt perfekt zu meinem Ziel, mich auf data storytelling und agentic infrastructure im großen Maßstab zu konzentrieren. Ich stehe ab sofort zur vollen Vollzeit zur Verfügung und bin für die Übergangsphase auch an einer Teilzeitregelung offen. Ich befinde mich derzeit in Deutschland und kann innerhalb von zwei bis drei Wochen innerhalb des Landes umziehen. Ich bin Blue Card-ready, was bedeutet, dass Ihr Legal Team lediglich den Vertrag ausstellen muss, ohne zusätzlichen Sponsorship-Aufwand.

\closing{Mit freundlichen Grüßen,}{Radheem Bin Razi}

\end{document}
